\documentclass[10pt]{article}
% BEGIN run values: written by code/python/weekly_report.py from weekly_simulation/results/state_t251_mock__doses_51x16x16_mock_15wk__2021/
\providecommand{\RunTag}{}\renewcommand{\RunTag}{code/\allowbreak{}python/\allowbreak{}state\_\allowbreak{}t251\_\allowbreak{}mock\_\allowbreak{}\_\allowbreak{}doses\_\allowbreak{}51x16x16\_\allowbreak{}mock\_\allowbreak{}15wk\_\allowbreak{}\_\allowbreak{}2021}
\providecommand{\RunInitialFile}{}\renewcommand{\RunInitialFile}{weekly\_\allowbreak{}simulation/\allowbreak{}input/\allowbreak{}state\_\allowbreak{}t251\_\allowbreak{}mock.csv}
\providecommand{\RunDosesFile}{}\renewcommand{\RunDosesFile}{weekly\_\allowbreak{}simulation/\allowbreak{}input/\allowbreak{}doses\_\allowbreak{}51x16x16\_\allowbreak{}mock\_\allowbreak{}15wk.xlsx}
\providecommand{\RunParams}{}\renewcommand{\RunParams}{2021}
\providecommand{\RunTvacc}{}\renewcommand{\RunTvacc}{347}
\providecommand{\RunTvaccDate}{}\renewcommand{\RunTvaccDate}{12 December 2020}
\providecommand{\RunWeeks}{}\renewcommand{\RunWeeks}{15}
\providecommand{\RunTlast}{}\renewcommand{\RunTlast}{445}
\providecommand{\RunTend}{}\renewcommand{\RunTend}{452}
\providecommand{\RunTendDate}{}\renewcommand{\RunTendDate}{27 March 2021}
\providecommand{\RunDaysNoVacc}{}\renewcommand{\RunDaysNoVacc}{96}
\providecommand{\RunFigDir}{}\renewcommand{\RunFigDir}{../weekly_simulation/results/state_t251_mock__doses_51x16x16_mock_15wk__2021/figures/}
\providecommand{\RunSwitch}{}\renewcommand{\RunSwitch}{$t=398$, $t=460$ and $t=548$ (1 February 2021, 4 April 2021, 1 July 2021)}
\providecommand{\RunRhoBetter}{}\renewcommand{\RunRhoBetter}{15}
\providecommand{\RunRhoRedMin}{}\renewcommand{\RunRhoRedMin}{8.7}
\providecommand{\RunRhoRedMax}{}\renewcommand{\RunRhoRedMax}{22}
\providecommand{\RunRhoRedMean}{}\renewcommand{\RunRhoRedMean}{13}
\providecommand{\RunRhoOptAbove}{}\renewcommand{\RunRhoOptAbove}{0}
\providecommand{\RunRhoUniAbove}{}\renewcommand{\RunRhoUniAbove}{0}
\providecommand{\RunCapWeeks}{}\renewcommand{\RunCapWeeks}{11}
\providecommand{\RunVPeak}{}\renewcommand{\RunVPeak}{61.1}
\providecommand{\RunVPeakDate}{}\renewcommand{\RunVPeakDate}{27 March 2021}
\providecommand{\RunInfPeak}{}\renewcommand{\RunInfPeak}{1.05}
\providecommand{\RunInfPeakDate}{}\renewcommand{\RunInfPeakDate}{12 December 2020}
\providecommand{\RunHospPeak}{}\renewcommand{\RunHospPeak}{0.0708}
\providecommand{\RunHospPeakDate}{}\renewcommand{\RunHospPeakDate}{12 December 2020}
\providecommand{\RunDeathsStart}{}\renewcommand{\RunDeathsStart}{60}
\providecommand{\RunDeathsEnd}{}\renewcommand{\RunDeathsEnd}{79}
\providecommand{\RunSeroPct}{}\renewcommand{\RunSeroPct}{6.6}
\providecommand{\RunDosesOne}{}\renewcommand{\RunDosesOne}{87}
\providecommand{\RunDosesTwo}{}\renewcommand{\RunDosesTwo}{52}
\providecommand{\RunDosesThree}{}\renewcommand{\RunDosesThree}{0}
\providecommand{\CmpEigenvectorDeathsEnd}{}\renewcommand{\CmpEigenvectorDeathsEnd}{78}
\providecommand{\CmpEigenvectorRhoAbove}{}\renewcommand{\CmpEigenvectorRhoAbove}{0}
\providecommand{\CmpEigenvectorInfPeak}{}\renewcommand{\CmpEigenvectorInfPeak}{1.05}
\providecommand{\CmpEigenvectorHospPeak}{}\renewcommand{\CmpEigenvectorHospPeak}{0.0708}
\providecommand{\CmpOptimalDeathsEnd}{}\renewcommand{\CmpOptimalDeathsEnd}{79}
\providecommand{\CmpOptimalRhoAbove}{}\renewcommand{\CmpOptimalRhoAbove}{0}
\providecommand{\CmpOptimalInfPeak}{}\renewcommand{\CmpOptimalInfPeak}{1.05}
\providecommand{\CmpOptimalHospPeak}{}\renewcommand{\CmpOptimalHospPeak}{0.0708}
% END run values
% the four figures, in the folder \RunFigDir of the run (relative to this file)
\newcommand{\figrho}{fig_weekly_optimal_rho.png}
\newcommand{\figcomp}{fig_weekly_optimal_compartments.png}
\newcommand{\figinf}{fig_weekly_optimal_infectious.png}
\newcommand{\figdeaths}{fig_weekly_optimal_deaths.png}
\usepackage[margin=0.65in]{geometry}
\usepackage{amsmath,amssymb}
\usepackage{graphicx}
\usepackage{xurl}
\usepackage[colorlinks=true,urlcolor=blue]{hyperref}
\pagestyle{empty}
% less space above a paragraph heading
\makeatletter
\renewcommand\paragraph{\@startsection{paragraph}{4}{\z@}{0.8ex \@plus 0.3ex}{-1em}%
  {\normalfont\normalsize\bfseries}}
\makeatother

\begin{document}
\begin{center}
{\large\bfseries The weekly allocation of vaccine: simulation, 2020--2022}\\[2pt]
Authors\footnote{Parts of this note were prepared with the assistance of an AI tool (Claude
Opus 5.5, Anthropic). The authors are solely responsible for all statements.}
\end{center}

\begin{center}
\includegraphics[width=0.36\textwidth]{\RunFigDir\figrho}\hspace{0.04\textwidth}
\includegraphics[width=0.36\textwidth]{\RunFigDir\figcomp}\\[2pt]
\includegraphics[width=0.36\textwidth]{\RunFigDir\figinf}\hspace{0.04\textwidth}
\includegraphics[width=0.36\textwidth]{\RunFigDir\figdeaths}
\end{center}

\paragraph{The simulation.}
The model of Feng, Day and Glasser (J.~Theoret.\ Biol.\ 633, 2026), 13 compartments
for each of $m=51$ locations (50 states and D.C.) and $n=16$ age groups, $10\,608$
unknowns, as in \texttt{DATA.pdf}, Section~2. The time $t$ is in days from 31 December
2019. With the parameter set \texttt{\RunParams} (\texttt{DATA.pdf},
Section~3), the system is solved by the Runge--Kutta method of order 8 (\texttt{DOP853})
from $t=251$ (7 September 2020), the initial time of Glasser et al.\ (J.~Theoret.\
Biol.\ 556, 2023, Section~3), to $t=\RunTend$ (\RunTendDate): with $\chi=0$ to $t=\RunTvacc$
(\RunTvaccDate), the start of vaccination of the set, then over \RunWeeks{} weeks,
$t=\RunTvacc,\dots,\RunTlast$, the number of weeks of the dose file. Each week the vaccination rates of the 3 doses are
those that minimize the spectral radius $\rho$ of the next generation matrix, subject to
the doses of the week of each location (all ages) and to at most $99\%$ of a group
vaccinated in a week; the minimization is started from the uniform allocation and from
the left eigenvector, and the smaller end value is kept. The values that change over time
change at \RunSwitch{} (\texttt{DATA.pdf}, Section~3.3). The program is \texttt{code/python/seir7\_ode.py
simulate optimal}; the figures are made by \texttt{code/python/plot\_weekly\_optimal.py}.

\paragraph{Results.}
In \RunRhoBetter{} of the \RunWeeks{} weeks the optimal allocation gives a smaller $\rho$ than
the uniform one, by $\RunRhoRedMin$--$\RunRhoRedMax\%$ ($\RunRhoRedMean\%$ on average);
$\rho>1$ in \RunRhoOptAbove{} weeks (optimal) and \RunRhoUniAbove{} weeks (uniform). In
\RunCapWeeks{} weeks the bound of $99\%$ a week is reached by some group. The vaccinated
$V$ peak at $\RunVPeak$ million (\RunVPeakDate); the infectious peak at $\RunInfPeak$
million (\RunInfPeakDate) and the hospitalized at $\RunHospPeak$ million
(\RunHospPeakDate). The deaths increase from $\RunDeathsStart$ thousand at $t=\RunTvacc$ to
$\RunDeathsEnd$ thousand at $t=\RunTend$.

\paragraph{Public data.}
Nine of the 11 inputs are built from published or public sources, each described in
\texttt{DATA.pdf}, Section~1, with its source and script: the population $N$ (Census
Vintage 2024), the contact matrix $U$ (Prem, Cook, Jit 2017), the distances and activity
factors (Census centers of population, county estimates, Gazetteer), the contact matrix
$C$ (from $U$, $L$, $d$), the stringency $s_i(t)$ (OxCGRT), the age-specific values
and $\beta^0$ (Glasser et al.\ 2023, Table~2 and Fig.~4), and the days of the variants (CDC
variant proportions). The constants are those of the set \texttt{\RunParams}
(\texttt{DATA.pdf}, Section~3). Rebuilding all 11 inputs from the files in
\texttt{data/Downloaded\_from\_source/} with the scripts named in \texttt{DATA.pdf} gives
files with contents identical to those used.

\paragraph{Mock data.}
The other two inputs are mock (\texttt{DATA.pdf}, Section~1.9), and the results above
depend on them.
(1)~The initial state at $t=251$, \texttt{\RunInitialFile}: from the public CDC
seroprevalence survey, rounds 3 and 4, by made-up rules (\texttt{DATA.pdf},
Section~1.9.1) in place of the procedure of Glasser et al.\ (2023, Section~3);
$\RunSeroPct\%$ seropositive.
(2)~The weekly doses by location, age group and dose,
\texttt{\RunDosesFile}: made up, \RunWeeks{} weeks, $\RunDosesOne$, $\RunDosesTwo$ and
$\RunDosesThree$ million first, second and third doses in total; to be replaced by more realistic data
later.


\begin{center}
{\large\bfseries Comparison with the eigenvector strategy}
\end{center}

\paragraph{The two strategies.}
The model, the inputs and the weeks are those above. Each week the doses of each location
and dose, summed over the ages, are the same for the two strategies; they differ only in how
these doses are
allocated to the age groups of the location (\texttt{DATA.pdf}, Section~2.3):
\begin{itemize}
\item \emph{optimal}: the rates that minimize the spectral radius $\rho$ of the next generation
matrix;
\item \emph{eigenvector}: the alternative strategy of Feng, Day and Glasser (J.~Theoret.\
Biol.\ 633, 2026, Section~2.2), the doses in proportion to the left eigenvector of the next
generation matrix at the end of the previous week, with the choices (C1)--(C4) of
\texttt{DATA.pdf}, Section~2.3.
\end{itemize}
Each strategy has its own solution of the system of ODEs from $t=\RunTvacc$; $\rho$ is
computed each week for the rates applied, with the state of that strategy. All results are
from \texttt{weekly\_\allowbreak{}simulation/\allowbreak{}results/\allowbreak{}\RunTag/}.

\paragraph{Results of the two strategies.}
\begin{center}
\begin{tabular}{lcc}
 & optimal & eigenvector\\\hline
weeks with $\rho>1$ & \CmpOptimalRhoAbove & \CmpEigenvectorRhoAbove\\
peak of the infectious (millions) & \CmpOptimalInfPeak & \CmpEigenvectorInfPeak\\
peak of the hospitalized (millions) & \CmpOptimalHospPeak & \CmpEigenvectorHospPeak\\
deaths at $t=\RunTend$ (thousands) & \CmpOptimalDeathsEnd & \CmpEigenvectorDeathsEnd
\end{tabular}
\end{center}

\begin{center}
\includegraphics[width=0.48\textwidth]{\RunFigDir fig_compare_rho.png}\hfill
\includegraphics[width=0.48\textwidth]{\RunFigDir fig_compare_infectious.png}\\[4pt]
\includegraphics[width=0.48\textwidth]{\RunFigDir fig_compare_hospitalized.png}\hfill
\includegraphics[width=0.48\textwidth]{\RunFigDir fig_compare_deaths.png}
\end{center}

\paragraph{Doses by age.} The doses of each type given to each age group, summed over the
weeks and the locations.
\begin{center}
\includegraphics[width=0.48\textwidth]{\RunFigDir fig_compare_doses_age_dose1.png}\hfill
\includegraphics[width=0.48\textwidth]{\RunFigDir fig_compare_doses_age_dose2.png}\\[4pt]
\includegraphics[width=0.48\textwidth]{\RunFigDir fig_compare_doses_age_dose3.png}
\end{center}

\end{document}
